% \RequirePackage{fontspec} % in unicode-math
\ifJHPreamble@omitUnicodeMath
    \RequirePackage{fontspec}
\else
    \RequirePackage[
        math-style=ISO, % Macht gr. Großbuchstaben kursiv
        bold-style=ISO 
        ]{unicode-math} 
\fi

\ifJHPreamble@PackagesByPath
    \newcommand{\__FontPath}{\JHPreamble@SetupPath/../../fonts/opentype/}
\else
    \newcommand{\__FontPath}{} % it will probably retry and find it in the System installation
\fi

\ifJHPreamble@StandardMathFonts
    % \setmathfont[
    %     version=sans
    % ]{} % There is no real standard sans math font compatible with unicode math
\else
    \ifJHPreamble@KpMath
        %%% kpfonts haben einige Zeichen nicht, deshalb zu STIX2Math gewechselt. Passen beide nicht wirklich zu MinionPro

        % \RequirePackage{kpfonts}
        % \RequirePackage[mathcal,bbsans,notext]{kpfonts-otf}
        \setmathfont[
            Path=\__FontPath,
            version=normal
        ]{KpMath-Regular.otf}

        \setmathfont[
            Path=\__FontPath,
            version=bold
        ]{KpMath-Bold.otf}

        % % \setmathfont{KpMath-Semibold.otf}[Path=\__FontPath,version=semibold]
        % \setmathfont{KpMath-Sans.otf}[Path=\__FontPath,version=sans]
        % \setmathfont{KpMath-Light.otf}[Path=\__FontPath,version=light]


        % \newfontfamily{\fallbackfont}{STIX2Math}[Path=\__FontPath,
        %     Extension = .otf
        % ]

        % \DeclareTextFontCommand{\usefallbackfont}{\fallbackfont}

    \else
        \ifJHPreamble@GyrePagella
            \setmathfont[
                Path=\__FontPath,
                version=normal
            ]{texgyrepagella-math.otf}

            \setmathfont[
                Path=\__FontPath,
                version=bold
            ]{KpMath-Bold.otf}
        \else
            \setmathfont[
                Path=\__FontPath,
                version=bold
            ]{XITSMath-Bold.otf}

            %%% STIX2Math und XITS (STIX)
            \ifJHPreamble@STIX
                \setmathfont[
                    Path=\__FontPath,
                    version=normal
                ]{STIX2Math.otf}
            \else
                \setmathfont[
                    Path=\__FontPath,
                    version=normal
                ]{XITSMath.otf} 
            \fi
        \fi
    \fi


    % \setmathfont[
    %     Path=\__FontPath,
    %     version=sans
    % ]{GFSNeohellenicMath.otf}
    
    % \setmathfont[
    %     Path=\__FontPath,
    %     range={"2983-"29FC},
    %     version=sans
    % ]{XITSMath.otf}
    
    % Was named KpMath-Sans.otf, which collides with a font name TeX Live's own kpfonts-otf
    % package registers via fonts/latex/kpfonts-otf/KpMath-Sans.fontspec. That file declares
    % BoldFont=KpMath-SansBold as a *default* font feature for anything named "KpMath-Sans" -
    % applied by fontspec purely from the name match, regardless of our Path= pointing elsewhere -
    % and it overrides an explicit BoldFont= given here. Since we didn't ship a bold sans cut,
    % that default lookup failed with "the font KpMath-SansBold cannot be found", even though we
    % never asked for a bold sans font ourselves. Renaming our copy to KpMath-Sans-Local.otf avoids
    % the collision; now that we DO have the real KpMath-SansBold.otf (the file kpfonts-otf expects,
    % downloaded separately since kpfonts-otf ships it as a system font, not bundled in that CTAN
    % package's own files), we load it explicitly as BoldFont below rather than relying on the
    % name-triggered default, which stays inert now that our filename no longer collides.
    \setmathfont[
        Path=\__FontPath,
        % range={up, it, bb, frak, cal, scr, sfup, sfit, bfup,bfit, bfcal},
        version=sans,
        BoldFont=KpMath-SansBold.otf
    ]{KpMath-Sans-Local.otf}
    % \setmathfont[
    %     Path=\__FontPath,
    %     range={it,up, bb},
    %     version=sans
    % ]{FiraMath-Regular.otf}
    %%%



    \mathversion{normal}
\fi

%%% Text Fonts

\ifJHPreamble@StandardTextFonts
\else
    \ifJHPreamble@GyrePagella
        \setmainfont{texgyrepagella}[
            Path=\__FontPath,
            Extension = .otf,
            UprightFont = *-regular,
            BoldFont = *-bold,
            ItalicFont = *-italic,
            BoldItalicFont = *-bolditalic, 
            Language = German
        ]
        \setsansfont{texgyreheros}[
            Path=\__FontPath,
            Extension = .otf,
            UprightFont = *-regular,
            BoldFont = *-bold,
            ItalicFont = *-italic,
            BoldItalicFont = *-bolditalic, 
            Language = German
        ]
        % decrease font size by 1pt (these fonts are smaller than Minion and Myriad)
        
    \else
        % \IfFontExistsTF{MinionPro}{
        \setmainfont{MinionPro}[
            Path=\__FontPath,
            Extension = .otf,
            UprightFont = *-Medium,
            BoldFont = *-Bold, %SemiBold
            ItalicFont = *-It, %
            BoldItalicFont = *-BoldIt, %SemiboldIt
            Language = German
        ] % Minion Pro und Myriad Pro sind anscheinend zur privaten Verwendung legal (aus der Installation von Adobe Acrobat Reader zu bekommen)
        % }{}

        %   \IfFontExistsTF{Myriad-Pro}{
        \setsansfont{Myriad-Pro}[ %SansSerif 
            Path=\__FontPath, 
            UprightFont = *-Regular.ttf,
            BoldFont = *-Bold.ttf, %SemiBold
            ItalicFont = *-It.otf, %
            Language = German,
            % Script = Default, % Seems to be on CustomDefault although that doesn't exist in the documentaion
            Contextuals = ResetAll
        ]
        % }{}

        %%% Braucht kompliziert konvertierte und installierte Myriad und Minion Type 1 Schriften mit Paketen. (Für pdflatex)

        % \RequirePackage{alphabeta}
        % \RequirePackage[Path=\__FontPath,lf]{MinionPro}
        % \RequirePackage[Path=\__FontPath,
        % 	scale=1,
        % 	lf,
        % 	sansmath]{MyriadPro}
    \fi
    
    \newfontfamily\codeFont{CascadiaCode}[
        Path=\__FontPath, 
        UprightFont = *-Regular.ttf,
        BoldFont = *-Bold.ttf
    ]
    
    
    % icomoon.io, boxy-svg.com
    \newfontfamily\JHSymbolFont[Path=\__FontPath]{divides_multipleOf.ttf}
    \newcommand\teilt{\mathbin{\text{\JHSymbolFont{d}}}}
    \newcommand\assoz{\mathbin{\text{\JHSymbolFont{a}}}}
    \newcommand\vielf{\mathbin{\text{\JHSymbolFont{m}}}}
\fi
        
%%%% FallbackFont
% Mit Kp-Math notwendig, weil die das nicht hatten.
% Diese Kommandos könnte man wahrscheinlich auch durch \setmathfont{STIXMath2}[Unicode-range] irgendwie weglassen.

% \RequirePackage{newunicodechar}
% \newunicodechar{Δ}{\increment} % Könnte man machen, aber unnötig. Es gibt das Zeichen ∆, das ist dasselbe wie \increment, aber verschieden von Δ (\Delta)

% \newcommand{\MakeCharAlwaysUseFont}[2][\fallbackfont]{\newunicodechar{#2}{#1{#2}}}
% \newcommand{\MakeCharAlwaysUseFontAndTextMode}[2][\fallbackfont]{\newunicodechar{#2}{\text{#1{#2}}}}

% \NewDocumentCommand\MakeCharactersAlwaysUseFallbackFont{ > { \SplitList {,} } m }{ \ProcessList{#1}{\MakeCharAlwaysUseFont} }
% \NewDocumentCommand\MakeCharactersAlwaysUseFallbackFontAndTextMode{ > { \SplitList {,} } m }{ \ProcessList{#1}{\MakeCharAlwaysUseFontAndTextMode} }


% \MakeCharactersAlwaysUseFallbackFontAndTextMode{𝟙}

% \MakeCharactersAlwaysUseFallbackFontAndTextMode{⥎,⥋,⥊,⥐}